\documentclass[10pt]{article}
\usepackage[a4paper,margin=0.6in]{geometry}
\usepackage{amsmath,amssymb,amsfonts,mathtools}
\usepackage[dvipsnames]{xcolor}
\usepackage{array,booktabs,float}
\usepackage[colorlinks=true,linkcolor=blue]{hyperref}
\allowdisplaybreaks
\numberwithin{equation}{section}
\setlength{\jot}{4pt}
\setlength{\emergencystretch}{3em}
\newcommand{\KK}{\mathcal K}
\setcounter{tocdepth}{1}

\title{Group I, Row V, part 2b:\\ the UV divergence in dimensional regularization ($\overline{\rm MS}$)}
\date{}

\begin{document}
\maketitle
\tableofcontents

%==========================================================================================
\section{Result}\label{sec:result}
Part 2b, as you wrote it, has one UV divergence. It is at $x\to z$, where the two gluons of the loop (at $x$ with $p^+$, at $z$ with $k^+-p^+$) come together. We add and subtract the short-distance limit of the integrand (the counterterm, Sec.~\ref{sec:ct}).
\begin{itemize}
\item The counterterm is integrated over $x$ in $d_\perp=2-\epsilon$ dimensions in $\overline{\rm MS}$, and over $p^+$ exactly. This gives the \textbf{UV term}.
\item What remains is the \textbf{rest}. It is UV finite and is evaluated in two dimensions.
\end{itemize}
Nothing is expanded in $\Lambda/k^+$ or $k^+/\vee$.
\begin{equation}
\boxed{\;\frac{d^{3}N_{2b}}{d^{3}k}=\frac{d^{3}N_{2b}}{d^{3}k}\bigg|_{\rm UV}+\frac{d^{3}N_{2b}}{d^{3}k}\bigg|_{\rm rest}+{\rm c.c.}\;}
\label{eq:split}
\end{equation}

\paragraph{The UV term (dimensional regularization, $d_\perp=2-\epsilon$, $\overline{\rm MS}$ scale $\mu$).}
\begin{align}
\frac{d^{3}N_{2b}}{d^{3}k}\bigg|_{\rm UV}={}&\frac{1}{(2\pi)^{3}}\,\frac{g^{4}N_{c}}{16\pi^{4}}\,\frac{1}{k^{+}}
\left[2\log\Big(\frac{\vee}{\Lambda}\Big)-2\log\Big(\frac{\vee}{k^{+}}\Big)+2\log\Big(1-\frac{\Lambda}{k^{+}}\Big)-\frac{11}{6}+\frac{4\Lambda}{k^{+}}-\frac{\Lambda^{2}}{k^{+2}}+\frac{2\Lambda^{3}}{3k^{+3}}\right]\notag\\
&\times\int_{y,z,x',w'}e^{-ik\cdot(w'-z)}\left[\frac{2}{\epsilon}+2\gamma_{E}+\log\frac{(y-z)^{2}\mu^{2}_{\overline{MS}}}{4}\right]
\frac{(x'-w')\cdot(y-z)}{(x'-w')^{2}\,(y-z)^{2}}\notag\\
&\times\Big[U^{ab'}(x')-U^{ab'}(w')\Big]\Big[U^{af}(y)-U^{af}(z)\Big]\rho^{b'}(x')\,\rho^{f}(y).
\label{eq:UV}
\end{align}
The bracket with the pole comes from the $x$ integral of the counterterm (Sec.~\ref{sec:dimreg}). In transverse coordinate space it is
\begin{equation}
\frac{2}{\epsilon}+2\gamma_{E}+\log\frac{(y-z)^{2}\mu^{2}_{\overline{MS}}}{4}
=\frac{1}{\pi}\int_{x}\Big[K(z,z,x)-K(y,z,x)\Big]
=\frac{1}{\pi}\int_{x}\frac{(y-x)\cdot(y-z)}{(y-x)^{2}\,(x-z)^{2}}\qquad(d_\perp=2-\epsilon),
\label{eq:coord}
\end{equation}
with $K(a,b,c)=\frac{(a-c)\cdot(b-c)}{(a-c)^2(b-c)^2}$. This is the same form as in your G1-II, $\frac{1}{\pi}\int_z[K(x,w,z)-K(w,w,z)]=-\big[\frac2\epsilon+2\gamma+\log\frac{(x-w)^2\mu^2_{\overline{MS}}}{4}\big]$.

\paragraph{After the $\overline{\rm MS}$ subtraction} (remove $\frac2\epsilon$; the $\gamma_E$ and $\log4\pi$ of $\overline{\rm MS}$ are already in $\mu_{\overline{MS}}$):
\begin{equation}
\frac{2}{\epsilon}+2\gamma_{E}+\log\frac{(y-z)^{2}\mu^{2}_{\overline{MS}}}{4}\;\longrightarrow\;\log\frac{(y-z)^{2}\mu^{2}_{\overline{MS}}\,e^{2\gamma_E}}{4}
\quad\text{in Eq.~\eqref{eq:UV}.}
\label{eq:msbar}
\end{equation}

\paragraph{The rest (UV finite; $d_\perp=2$; no plus prescription).}
\begin{align}
\frac{d^{3}N_{2b}}{d^{3}k}\bigg|_{\rm rest}={}&\frac{1}{(2\pi)^{3}}\,\frac{g^{4}}{8\pi^{4}}\,\frac{1}{k^{+}}\int_{\Lambda}^{k^{+}-\Lambda}\frac{dp^{+}}{2\pi}\int_{x,y,z,x',w'}e^{-ik\cdot(w'-z)}\,\Bigg\{\notag\\
&e^{-ik\cdot(z-x)\frac{p^{+}}{k^{+}}}\,\frac{1}{\left[p^{+}(y-x)^{2}+(k^{+}-p^{+})(y-z)^{2}\right]}\,\frac{(x'-w')^{i}\,(z-x)^{m}}{(x'-w')^{2}\,(z-x)^{2}}\notag\\
&\times\Bigg[\frac{\delta^{im}\,d_{\perp}}{2\,(x-z)^{2}}\left[(y-x)^{2}-(y-z)^{2}\right]\frac{p^{+}(k^{+}-p^{+})}{k^{+2}}
-\frac{p^{+}}{k^{+}-p^{+}}\left[\frac{(y-x)^{i}(x-z)^{m}}{(x-z)^{2}}-\frac{(y-x)^{i}(y-z)^{m}}{2(y-z)^{2}}\right]\notag\\
&\qquad+\frac{(k^{+}-p^{+})\,\delta_{im}}{k^{+}}\left[\frac{(y-z)\cdot(x-z)}{(x-z)^{2}}+\frac{(y-x)\cdot(y-z)}{2(y-x)^{2}}-\frac{(y-x)^{2}-(y-z)^{2}}{2(x-z)^{2}}\right]\notag\\
&\qquad-\frac{k^{+}-p^{+}}{p^{+}}\left[\frac{(y-z)^{i}(x-z)^{m}}{(x-z)^{2}}+\frac{(y-x)^{m}(y-z)^{i}}{2(y-x)^{2}}\right]\notag\\
&\qquad+\frac{p^{+}\,\delta_{im}}{k^{+}}\left[\frac{(y-x)\cdot(x-z)}{(x-z)^{2}}-\frac{(y-x)\cdot(y-z)}{2(y-z)^{2}}-\frac{(y-x)^{2}-(y-z)^{2}}{2(x-z)^{2}}\right]\notag\\
&\qquad-\left[\frac{(y-x)^{m}(x-z)^{i}}{(x-z)^{2}}-\frac{(y-x)^{m}(y-z)^{i}}{2(y-z)^{2}}+\frac{(y-z)^{m}(x-z)^{i}}{(x-z)^{2}}+\frac{(y-x)^{i}(y-z)^{m}}{2(y-x)^{2}}\right]\Bigg]\notag\\
&\times\Big[U^{ab'}(x')-U^{ab'}(w')\Big]\Big[N_{c}\,U^{af}(y)-N_{c}\,U^{af}(x)\Big]\rho^{b'}(x')\,\rho^{f}(y)\notag\\
&-\frac{1}{k^{+}}\left[\frac{p^{+}(k^{+}-p^{+})}{k^{+2}}+\frac{p^{+}}{k^{+}-p^{+}}+\frac{k^{+}-p^{+}}{p^{+}}\right]
\frac{(x'-w')\cdot(y-z)}{(x'-w')^{2}\,(y-z)^{2}}\,\frac{(y-x)\cdot(y-z)}{(y-x)^{2}\,(x-z)^{2}}\notag\\
&\times\Big[U^{ab'}(x')-U^{ab'}(w')\Big]\Big[N_{c}\,U^{af}(y)-N_{c}\,U^{af}(z)\Big]\rho^{b'}(x')\,\rho^{f}(y)\Bigg\}.
\label{eq:rest}
\end{align}
Everything before the last two lines is your expression for 2b, unchanged (in the rest $d_\perp=2$). The last two lines subtract the counterterm.

\paragraph{Notes.}
\begin{enumerate}
\item \textbf{Conventions.} $d_\perp=2-\epsilon$, as in your G1-II, so $\frac2\epsilon=\frac{1}{\epsilon_{\rm std}}$ in the usual $D=4-2\epsilon_{\rm std}$ notation. Every $d_\perp$-dimensional transverse integral carries $\big(\frac{\mu^2e^{\gamma_E}}{4\pi}\big)^{\epsilon/2}$, $\mu\equiv\mu_{\overline{MS}}$. In $\overline{\rm MS}$ only $\frac2\epsilon$ is then subtracted.
\item \textbf{The log rule.} The $p^+$ integral gives $2\log\frac{k^+-\Lambda}{\Lambda}$. We used $\log\frac{k^+}{\Lambda}=\log\frac\vee\Lambda-\log\frac{\vee}{k^+}$, i.e.\ your $\log\frac{\Lambda}{k^+}=-\log\frac\vee\Lambda+\log\frac\vee{k^+}$, so $\log\frac{k^+-\Lambda}{\Lambda}=\log\frac\vee\Lambda-\log\frac\vee{k^+}+\log\big(1-\frac{\Lambda}{k^+}\big)$.
\item \textbf{The UV term is LO times a factor.} Its prefactor is $\frac{\alpha_sN_c}{2\pi}$ times that of LO. The integral is the LO cross section with $x\to y$ (the charge) and $w\to z$ (the gluon), with the bracket of Eq.~\eqref{eq:coord} inside the integral.
\item \textbf{Running coupling.} The $-\frac{11}{6}$ term is $-\frac{\alpha_s\beta_0}{4\pi}\big[\frac2\epsilon+2\gamma_E+\log\frac{(y-z)^2\mu^2}{4}\big]\times$LO, with $\beta_0=\frac{11}{3}N_c$. Your G1-II (part 1) has the same term with the opposite sign, $+\frac{\alpha_s\beta_0}{4\pi}\big[\frac2\epsilon+2\gamma_E+\log\frac{(x-w)^2\mu^2}{4}\big]\times$LO. After renaming $y\to x$, $z\to w$, the two cancel exactly per ordering, pole and logarithm included (Sec.~\ref{sec:beta}).
\item \textbf{The pole multiplying the rapidity logarithms} is not a renormalization of the coupling. In the form of Eq.~\eqref{eq:coord}, it is an $x$ integral with a $1/(x-z)^2$ singularity, the same type as your G1-II after F6. These terms are part of the $\log(\vee/\Lambda)$ two-$\rho$ terms. Their sum is JIMWLK$\otimes$LO, where such singularities cancel between rows (Sec.~\ref{sec:rap}).
\item \textbf{Choice of counterterm.} The counterterm is cut off at the distance $|y-z|$ through $K(z,z,x)-K(y,z,x)$. The cut $\theta(|y-z|-|x-z|)/(x-z)^2$ and the smooth form $\frac{1}{(x-z)^2}-\frac{1}{(x-z)^2+(y-z)^2}$ give the same UV term. They also give the same rest after the $x$ integration (Sec.~\ref{sec:ct}).
\end{enumerate}

%==========================================================================================
\section{Notation and scheme}\label{sec:conv}
\begin{itemize}
\item $\xi\equiv p^+/k^+$, $\mathbf r\equiv x-z$, $Y\equiv y-z$; $\KK(X)\equiv X/X^2$.
\item $K(a,b,c)\equiv\frac{(a-c)\cdot(b-c)}{(a-c)^2(b-c)^2}$, as in your Appendix E: the third argument is the point that is integrated over. So $K(z,z,x)=\frac{1}{(x-z)^2}$, and $K(y,z,x)=\frac{(y-x)\cdot(z-x)}{(y-x)^2(z-x)^2}$.
\item The phase $e^{-ik\cdot(z-x)p^+/k^+}$ of your expression depends on $x$, $z$ and $p^+$, so it stands inside $\int_{x,z}$ and $\int dp^+$. We write all transverse integrals together, $\int_{x,y,z,x',w'}$.
\item \textbf{Scheme} (the same as for G1-II):
\begin{itemize}
\item the loop integral, over the size $x-z$ of the gluon pair, is done in $d_\perp=d=2-\epsilon$ at fixed $y,z,x',w'$;
\item the gluon polarizations are $d$-dimensional (your $d_\perp$ in the first term of the bracket);
\item the external factors stay two-dimensional.
\end{itemize}
\item $\alpha_s=\frac{g^2}{4\pi}$; $\beta_0=\frac{11}{3}N_c$ (gluons only).
\item LO (\texttt{main.tex}):
\begin{equation*}
\frac{d^3N_{\rm LO}}{d^3k}=\frac{1}{(2\pi)^3}\frac{g^2}{2\pi^2k^+}\int_{x,w,x',w'}e^{-ik\cdot(w'-w)}\frac{(x-w)\cdot(x'-w')}{(x-w)^2(x'-w')^2}[U(x)-U(w)]^{ab}[U(x')-U(w')]^{ab'}\rho^b(x)\rho^{b'}(x').
\end{equation*}
\end{itemize}

%==========================================================================================
\section{Where part 2b is UV divergent}\label{sec:uv}
\subsection{The regions}
\begin{itemize}
\item $x\to z$: the integrand behaves like $1/(x-z)^2$, and the colour factor $N_c[U(y)-U(x)]\to N_c[U(y)-U(z)]$ does not vanish. $\int d^2x/(x-z)^2$ diverges logarithmically. This is the UV divergence.
\item $y\to x$ and $y\to z$: at most $1/|y-x|$ or $1/|y-z|$, which is integrable.
\item $x,y,z$ all together: the integrand scales as $\rho^{-3}$ against $d^2x\,d^2y\sim\rho^3d\rho$, and the colour factor vanishes as well. Integrable.
\item Large $|x-z|$: the integrand falls like $1/(x-z)^2$. Its angular average is not zero, so the $x$ integral converges because of the phase $e^{-ik\cdot(z-x)p^+/k^+}$ (checked numerically, Table~\ref{tab:checks}). This holds for your expression as written, and the subtraction below does not change it: the counterterm falls like $|x-z|^{-3}$.
\end{itemize}

\subsection{The short-distance limit}\label{sec:short}
Put $z$ at the origin, $x=\mathbf r$, $y=Y$, and let $r\to0$. Then:
\begin{itemize}
\item $e^{-ik\cdot(w'-z)}e^{-ik\cdot(z-x)\xi}\to e^{-ik\cdot(w'-z)}$;
\item $\frac{1}{p^+(y-x)^2+(k^+-p^+)(y-z)^2}\to\frac{1}{k^+(y-z)^2}$;
\item $N_c[U(y)-U(x)]\to N_c[U(y)-U(z)]$.
\end{itemize}
Call the six terms of your bracket $B_1,\dots,B_6$ in the order written, without their $p^+$ factors. Contracted with $\frac{(z-x)^m}{(z-x)^2}=-\frac{r^m}{r^2}$, their leading terms are:
\begin{center}\small
\begin{tabular}{@{}c c l l@{}}
term & $p^+$ factor & leading term of $-\frac{r^m}{r^2}B^{im}$ & angular average in $d$ dimensions\\ \hline
$B_1$ & $\xi(1-\xi)$ & $d_\perp\,\frac{r^i(Y\cdot r)}{r^4}-\frac{d_\perp}{2}\frac{r^i}{r^2}$ & $\frac{d_\perp}{d}\frac{Y^i}{r^2}=\frac{Y^i}{r^2}$\\
$B_2$ & $-\frac{\xi}{1-\xi}$ & $-\frac{Y^i}{r^2}+\frac{r^i}{r^2}+\frac{(Y\cdot r)Y^i}{2Y^2r^2}$ & $-\frac{Y^i}{r^2}$\\
$B_3$ & $1-\xi$ & $-2\frac{r^i(Y\cdot r)}{r^4}$ & $-\frac2d\frac{Y^i}{r^2}$\\
$B_4$ & $-\frac{1-\xi}{\xi}$ & $-\frac{Y^i}{r^2}-\frac{(Y\cdot r)Y^i}{2Y^2r^2}$ & $-\frac{Y^i}{r^2}$\\
$B_5$ & $\xi$ & $-2\frac{r^i(Y\cdot r)}{r^4}+2\frac{r^i}{r^2}$ & $-\frac2d\frac{Y^i}{r^2}$\\
$B_6$ & $-1$ & $-2\frac{r^i(Y\cdot r)}{r^4}+\frac{r^i}{r^2}$ & $-\frac2d\frac{Y^i}{r^2}$\\
\end{tabular}
\end{center}
We used $\langle r^ir^j\rangle=\delta^{ij}r^2/d$; terms odd in $\mathbf r$ average to zero. Multiplying by the $p^+$ factors and adding,
\begin{equation}
\xi(1-\xi)+\frac{\xi}{1-\xi}-\frac2d(1-\xi)+\frac{1-\xi}{\xi}-\frac2d\,\xi+\frac2d
=\xi(1-\xi)+\frac{\xi}{1-\xi}+\frac{1-\xi}{\xi}.
\end{equation}
The result does not depend on $d$:
\begin{itemize}
\item the $d_\perp$ of the polarization sum in $B_1$ cancels the $1/d$ of the angular average;
\item the $1/d$ terms of $B_3$, $B_5$, $B_6$ cancel among themselves.
\end{itemize}
So no $O(\epsilon)$ term of the coefficient multiplies the pole. With $\frac{p^+(k^+-p^+)}{k^{+2}}+\frac{p^+}{k^+-p^+}+\frac{k^+-p^+}{p^+}=\xi(1-\xi)+\frac{\xi}{1-\xi}+\frac{1-\xi}{\xi}$, the integrand of 2b at $x\to z$ is
\begin{align}
&\frac{1}{(2\pi)^{3}}\frac{g^{4}}{8\pi^{4}}\frac{1}{k^{+}}\frac{dp^+}{2\pi}\,e^{-ik\cdot(w'-z)}\,\frac{1}{k^{+}}\left[\frac{p^{+}(k^{+}-p^{+})}{k^{+2}}+\frac{p^{+}}{k^{+}-p^{+}}+\frac{k^{+}-p^{+}}{p^{+}}\right]
\frac{(x'-w')\cdot(y-z)}{(x'-w')^{2}(y-z)^{2}}\,\frac{1}{(x-z)^{2}}\notag\\
&\times\Big[U^{ab'}(x')-U^{ab'}(w')\Big]\Big[N_cU^{af}(y)-N_cU^{af}(z)\Big]\rho^{b'}(x')\rho^{f}(y)
\;+\;\big(\text{zero angular average}\big)+O\Big(\frac{1}{|x-z|}\Big).
\label{eq:limit}
\end{align}
The bracket is the real part of the $g\to gg$ splitting function divided by $2N_c$. In $d$ dimensions every factor $\frac{r^m}{r^2}$ of the wave functions becomes $c_d\frac{r^m}{(r^2)^{d/2}}$ with $c_d=\Gamma(\frac d2)\pi^{1-d/2}$ ($c_2=1$), so $\frac{1}{(x-z)^2}\to c_d^2\,[(x-z)^2]^{1-d}$ with the same coefficient.

%==========================================================================================
\section{The counterterm}\label{sec:ct}
We subtract Eq.~\eqref{eq:limit} with $\frac{1}{(x-z)^2}$ replaced by
\begin{equation}
K(z,z,x)-K(y,z,x)=\frac{1}{(x-z)^2}-\frac{(y-x)\cdot(z-x)}{(y-x)^2(z-x)^2}=\frac{(y-x)\cdot(y-z)}{(y-x)^{2}\,(x-z)^{2}}.
\label{eq:S}
\end{equation}
The counterterm is therefore
\begin{align}
{\rm CT}={}&\frac{1}{(2\pi)^{3}}\frac{g^{4}}{8\pi^{4}}\frac{1}{k^{+}}\int_{\Lambda}^{k^+-\Lambda}\frac{dp^+}{2\pi}\int_{x,y,z,x',w'}e^{-ik\cdot(w'-z)}\,\frac{1}{k^{+}}\left[\frac{p^{+}(k^{+}-p^{+})}{k^{+2}}+\frac{p^{+}}{k^{+}-p^{+}}+\frac{k^{+}-p^{+}}{p^{+}}\right]\notag\\
&\times\frac{(x'-w')\cdot(y-z)}{(x'-w')^{2}(y-z)^{2}}\,\big[K(z,z,x)-K(y,z,x)\big]\Big[U^{ab'}(x')-U^{ab'}(w')\Big]\Big[N_cU^{af}(y)-N_cU^{af}(z)\Big]\rho^{b'}(x')\rho^{f}(y).
\label{eq:CT}
\end{align}
Then $\frac{d^3N_{2b}}{d^3k}=({\rm 2b}-{\rm CT})+{\rm CT}$. The first bracket is the rest, Eq.~\eqref{eq:rest}, and CT is the UV term, Eq.~\eqref{eq:UV}. The kernel \eqref{eq:S} has the following properties:
\begin{enumerate}
\item[(i)] \textbf{Same UV limit.} At $x\to z$ it is $\frac{1}{(x-z)^2}+\frac{(y-z)\cdot(x-z)}{(y-z)^2(x-z)^2}+O(1)$, and the second term has zero angular average. So the rest has no $1/(x-z)^2$ term left, for every $p^+$ in $\Lambda<p^+<k^+-\Lambda$.
\item[(ii)] \textbf{No new IR problem.} At large $|x-z|$ it falls like $|x-z|^{-3}$.
\item[(iii)] At $x\to y$ it behaves like $1/|x-y|$, which is integrable.
\item[(iv)] \textbf{It is a cut at $|x-z|=|y-z|$.} In two dimensions its average over the directions of $x-z$ at fixed $|x-z|$ is exactly
\begin{equation}
\Big\langle\frac{(y-x)\cdot(y-z)}{(y-x)^{2}(x-z)^{2}}\Big\rangle_{\rm angle}=\frac{\theta\big(|y-z|-|x-z|\big)}{(x-z)^{2}}.
\label{eq:angle}
\end{equation}
Proof: in complex notation $\frac{(y-x)}{(y-x)^2}=\frac{1}{\overline{Y}-\bar r}$. For $r<|Y|$, expanding in $\bar r/\overline Y$ gives the average $1/\overline Y$, i.e.\ $Y/Y^2$. For $r>|Y|$, expanding in $\overline Y/\bar r$ gives zero.
\end{enumerate}
By (iv), the cut $\frac{\theta(|y-z|-|x-z|)}{(x-z)^2}$, which I gave you before, has the same $x$ integral at every $y,z,p^+$. So does the smooth form $\frac{1}{(x-z)^2}-\frac{1}{(x-z)^2+(y-z)^2}$: in two dimensions both are $\pi\log\frac{(y-z)^2}{r_0^2}$ with a cutoff $r_0\to0$, and in $d$ dimensions they agree with Eq.~\eqref{eq:coord} (Table~\ref{tab:checks}). All three choices therefore give the same UV term, and the same rest after the $x$ integral. We use Eq.~\eqref{eq:S} because it is built from the kernels $K$ of your notes, and because it puts the UV term in the same form as your G1-II after F6.

%==========================================================================================
\section{The \texorpdfstring{$x$}{x} integral in \texorpdfstring{$d_\perp=2-\epsilon$}{d=2-eps} dimensions}\label{sec:dimreg}
In $d$ dimensions the kernels of Eq.~\eqref{eq:S} become
\begin{equation}
K_d(a,b,x)=c_d^2\,\frac{(a-x)\cdot(b-x)}{[(a-x)^2]^{d/2}\,[(b-x)^2]^{d/2}},\qquad c_d=\Gamma\Big(\frac d2\Big)\pi^{1-d/2},
\end{equation}
because
\begin{equation}
\int\frac{d^dq}{(2\pi)^d}\,e^{iq\cdot r}\,\frac{q^m}{q^2}=\frac{i\,\Gamma(d/2)}{2\pi^{d/2}}\frac{r^m}{(r^2)^{d/2}},\qquad\text{i.e.}\qquad c_d\frac{r^m}{(r^2)^{d/2}}=-2\pi i\int\frac{d^dq}{(2\pi)^d}\,e^{iq\cdot r}\,\frac{q^m}{q^2}.
\end{equation}
\paragraph{Step 1: $K_d(z,z,x)$.} It equals $c_d^2[(x-z)^2]^{1-d}$, which has no scale, so $\int d^dx\,K_d(z,z,x)=0$. In this integral the UV pole cancels against an IR pole. The IR pole cancels in turn against that of $K_d(y,z,x)$ in step 2, since the difference \eqref{eq:S} is IR finite.
\paragraph{Step 2: $K_d(y,z,x)$.} With the Fourier representation of both factors, the $x$ integral gives $(2\pi)^d\delta^{(d)}(q+q')$:
\begin{equation}
\int d^dx\,K_d(y,z,x)=(-2\pi i)^2\int\frac{d^dq}{(2\pi)^d}\,e^{iq\cdot(y-z)}\,\frac{q\cdot(-q)}{q^2\,q^2}=4\pi^2\int\frac{d^dq}{(2\pi)^d}\,\frac{e^{iq\cdot(y-z)}}{q^2}=\pi^{2-d/2}\,\Gamma\Big(\frac d2-1\Big)\big[(y-z)^2\big]^{1-d/2}.
\end{equation}
(For $d=3$ this is $\pi/|y-z|$, i.e.\ the known $\int d^3x\,\nabla\frac1{|x-y|}\cdot\nabla\frac1{|x-z|}=\frac{4\pi}{|y-z|}$.)
\paragraph{Step 3: $\overline{\rm MS}$.} With $d=2-\epsilon$ and the factor $\big(\frac{\mu^2e^{\gamma_E}}{4\pi}\big)^{\epsilon/2}$,
\begin{align}
\Big(\frac{\mu^2e^{\gamma_E}}{4\pi}\Big)^{\frac\epsilon2}\int d^dx\,\big[K_d(z,z,x)-K_d(y,z,x)\big]
&=-\pi\,\Gamma\Big(-\frac\epsilon2\Big)\Big(\frac{(y-z)^2\mu^2e^{\gamma_E}}{4}\Big)^{\frac\epsilon2}\notag\\
&=\pi\left[\frac2\epsilon+2\gamma_E+\log\frac{(y-z)^2\mu^2}{4}\right]+O(\epsilon),
\label{eq:dimreg}
\end{align}
using $\Gamma(-\frac\epsilon2)=-\frac2\epsilon-\gamma_E+O(\epsilon)$. This is Eq.~\eqref{eq:coord}.
\paragraph{Check with the cut.} Split $K_d(z,z,x)-K_d(y,z,x)=\theta(R-r)\,c_d^2(r^2)^{1-d}+\big[K_d(z,z,x)-K_d(y,z,x)-\theta(R-r)\,c_d^2(r^2)^{1-d}\big]$, with $R=|y-z|$.
\begin{itemize}
\item The first part gives
\begin{equation}
\Big(\frac{\mu^2e^{\gamma_E}}{4\pi}\Big)^{\frac\epsilon2}c_d^2\,\Omega_d\int_0^Rdr\,r^{d-1}r^{2-2d}=\frac{2\pi}{\epsilon}\,\Gamma\Big(1-\frac\epsilon2\Big)\Big(\frac{R^2\mu^2e^{\gamma_E}}{4}\Big)^{\frac\epsilon2}=\pi\left[\frac2\epsilon+2\gamma_E+\log\frac{R^2\mu^2}{4}\right]+O(\epsilon),
\end{equation}
with $\Omega_d=2\pi^{d/2}/\Gamma(d/2)$.
\item The second part is UV and IR finite, so it can be done at $d=2$, where it vanishes by Eq.~\eqref{eq:angle}.
\end{itemize}
The two results agree. For the smooth form, $\int_0^\infty dt\,\frac{t^{1-d}}{1+t^2}=\frac{\pi}{2\sin(\pi\epsilon/2)}=\frac1\epsilon+O(\epsilon)$ replaces $\int_0^1dt\,t^{1-d}=\frac1\epsilon$, so it also gives the same result.

%==========================================================================================
\section{The \texorpdfstring{$p^+$}{p+} integral (exact)}\label{sec:pplus}
With $\lambda=\Lambda/k^+$:
\begin{align}
\int_{\Lambda}^{k^+-\Lambda}\frac{dp^+}{k^+}\left[\frac{p^{+}(k^{+}-p^{+})}{k^{+2}}+\frac{p^{+}}{k^{+}-p^{+}}+\frac{k^{+}-p^{+}}{p^{+}}\right]
&=\int_\lambda^{1-\lambda}d\xi\Big[\xi(1-\xi)+\frac{1}{1-\xi}+\frac1\xi-2\Big]\notag\\
&=\Big[\frac16-\lambda^2+\frac23\lambda^3\Big]+2\log\frac{1-\lambda}{\lambda}-2(1-2\lambda)\notag\\
&=2\log\frac{k^+-\Lambda}{\Lambda}-\frac{11}{6}+\frac{4\Lambda}{k^+}-\frac{\Lambda^2}{k^{+2}}+\frac{2\Lambda^3}{3k^{+3}}.
\label{eq:xi}
\end{align}
This is exact; no term is dropped. With the log rule of Sec.~\ref{sec:result},
\begin{equation}
2\log\frac{k^+-\Lambda}{\Lambda}=2\log\frac{\vee}{\Lambda}-2\log\frac{\vee}{k^+}+2\log\Big(1-\frac{\Lambda}{k^+}\Big).
\end{equation}
The $-\frac{11}{6}$ comes from $\int_0^1[\xi(1-\xi)-2]\,d\xi$. The $\log$ comes from the $\frac1\xi$ and $\frac1{1-\xi}$ parts.

%==========================================================================================
\section{The UV term}\label{sec:UVterm}
\subsection{Assembling}
In Eq.~\eqref{eq:CT} we do the $x$ integral with Eq.~\eqref{eq:dimreg} and the $p^+$ integral with Eq.~\eqref{eq:xi}. Using $\frac{dp^+}{2\pi}\frac{1}{k^+}=\frac{1}{2\pi}\frac{dp^+}{k^+}$, the prefactor is
\begin{equation}
\frac{1}{(2\pi)^3}\,\frac{g^4}{8\pi^4}\,\frac1{k^+}\times\frac{1}{2\pi}\times\pi\times N_c=\frac{1}{(2\pi)^3}\,\frac{g^4N_c}{16\pi^4}\,\frac{1}{k^+}.
\end{equation}
This gives Eq.~\eqref{eq:UV}. Relative to the LO prefactor $\frac{1}{(2\pi)^3}\frac{g^2}{2\pi^2k^+}$ the factor is $\frac{g^2N_c}{8\pi^2}=\frac{\alpha_sN_c}{2\pi}$. So, with $L\equiv\big[2\log\frac\vee\Lambda-2\log\frac\vee{k^+}+2\log(1-\frac\Lambda{k^+})+\frac{4\Lambda}{k^+}-\frac{\Lambda^2}{k^{+2}}+\frac{2\Lambda^3}{3k^{+3}}\big]$,
\begin{equation}
\frac{d^{3}N_{2b}}{d^{3}k}\bigg|_{\rm UV}=\frac{\alpha_sN_c}{2\pi}\Big[L-\frac{11}{6}\Big]\times\frac{d^3N_{\rm LO}}{d^3k}\bigg|_{x\to y,\;w\to z}\Big[\text{with }\frac2\epsilon+2\gamma_E+\log\frac{(y-z)^2\mu^2}{4}\text{ inside}\Big].
\end{equation}
The two charges are in the order of your expression, $\rho^{b'}(x')\rho^f(y)$.

\subsection{\texorpdfstring{$\overline{\rm MS}$}{MS-bar}}
In $\overline{\rm MS}$ the pole $\frac2\epsilon=\frac1{\epsilon_{\rm std}}$ is subtracted. The $\gamma_E$ and $\log4\pi$ are already absorbed in $\mu_{\overline{MS}}$ through the factor $(\mu^2e^{\gamma_E}/4\pi)^{\epsilon/2}$. This gives Eq.~\eqref{eq:msbar}:
\begin{equation}
\frac{2}{\epsilon}+2\gamma_{E}+\log\frac{(y-z)^{2}\mu^{2}_{\overline{MS}}}{4}\;\to\;\log\frac{(y-z)^{2}\mu^{2}_{\overline{MS}}}{c_0},\qquad c_0=4e^{-2\gamma_E}.
\end{equation}

\subsection{The \texorpdfstring{$\beta_0$}{beta0} part and G1-II}\label{sec:beta}
The $-\frac{11}{6}$ term is
\begin{equation}
\frac{\alpha_sN_c}{2\pi}\Big(-\frac{11}{6}\Big)=-\frac{\alpha_s\beta_0}{4\pi},\qquad
\frac{d^{3}N_{2b}}{d^{3}k}\bigg|_{\rm UV,\,\beta_0}=-\frac{\alpha_s\beta_0}{4\pi}\Big[\frac2\epsilon+2\gamma_E+\log\frac{(y-z)^2\mu^2}{4}\Big]\times{\rm LO}(y,z).
\end{equation}
Your G1-II (part 1) has
\begin{equation*}
\frac{1}{(2\pi)^{3}}\frac{g^{4}N_{c}}{32\pi^{4}}\frac{1}{k^{+}}\int e^{-ik(w'-w)}\KK(x'-w')\cdot\KK(x-w)\Big[-\frac{2}{\epsilon}-2\gamma-\log\frac{(x-w)^{2}\mu^{2}}{4}\Big]\Big[\frac{11}{3}+\dots\Big][U(x')-U(w')][U(w)-U(x)]\rho\rho,
\end{equation*}
with the opposite colour sign, $[U(w)-U(x)]=-[U(x)-U(w)]$. Relative to LO its prefactor is $\frac{\alpha_sN_c}{4\pi}$. Per ordering, both rows multiply $\frac{\alpha_sN_c}{4\pi}\big[\frac2\epsilon+2\gamma_E+\log\frac{L_\perp^2\mu^2}{4}\big]\times$LO, with $L_\perp=x-w$ in G1-II and $L_\perp=y-z$ in 2b (the same after renaming). The coefficients are:
\begin{center}
\begin{tabular}{@{}l l@{}}
row & coefficient of $\frac{\alpha_sN_c}{4\pi}\big[\frac2\epsilon+2\gamma_E+\log\frac{L_\perp^2\mu^2}{4}\big]\times{\rm LO}$\\ \hline
G1-II (part 1) & $+\frac{11}{3}-2\log\frac\vee\Lambda+2\log\frac\vee{k^+}+\log\big(1-\frac{k^+}\vee\big)$\\
G1-V, 2b & $-\frac{11}{3}+4\log\frac\vee\Lambda-4\log\frac\vee{k^+}+4\log\big(1-\frac\Lambda{k^+}\big)+\frac{8\Lambda}{k^+}-\frac{2\Lambda^2}{k^{+2}}+\frac{4\Lambda^3}{3k^{+3}}$\\
sum & $2\log\frac\vee\Lambda-2\log\frac\vee{k^+}+\log\big(1-\frac{k^+}\vee\big)+4\log\big(1-\frac\Lambda{k^+}\big)+\frac{8\Lambda}{k^+}-\frac{2\Lambda^2}{k^{+2}}+\frac{4\Lambda^3}{3k^{+3}}$\\
\end{tabular}
\end{center}
For G1-II we used $\log(\frac\vee{k^+}-1)=\log\frac\vee{k^+}+\log(1-\frac{k^+}\vee)$. The $\beta_0$ terms cancel per ordering, including the pole and the $\mu$ dependence. So the sum G1-II + 2b needs no coupling renormalization at this order. As noted in \texttt{RowV\_UV.pdf} (Sec.~6.3), the $\frac{\alpha_s\beta_0}{4\pi}\frac1{\epsilon_{\rm std}}$ that the bare coupling of LO requires must then come from another row (the self-energy of the measured gluon after the target); this is still open.

\subsection{The pole multiplying the rapidity logarithms}\label{sec:rap}
The remaining poles multiply $\log\frac\vee\Lambda$, $\log\frac\vee{k^+}$, $\log(1-\frac\Lambda{k^+})$ and the powers of $\Lambda/k^+$. They come from the $\frac1\xi$ and $\frac{1}{1-\xi}$ parts of the splitting function, so they are not a renormalization of the coupling.
\begin{itemize}
\item In coordinate space, Eq.~\eqref{eq:coord}, the $\log\frac\vee\Lambda$ term is
\begin{align*}
&\frac{1}{(2\pi)^3}\frac{g^4N_c}{16\pi^5}\frac{1}{k^+}\,2\log\frac\vee\Lambda\int_{x,y,z,x',w'}e^{-ik\cdot(w'-z)}\frac{(x'-w')\cdot(y-z)}{(x'-w')^2(y-z)^2}\,\frac{(y-x)\cdot(y-z)}{(y-x)^2(x-z)^2}\\
&\qquad\times[U(x')-U(w')]^{ab'}[U(y)-U(z)]^{af}\rho^{b'}(x')\rho^f(y).
\end{align*}
This is a two-$\rho$ term with a $1/(x-z)^2$ singularity, like your G1-II after F6.
\item The sum of all $\log\frac\vee\Lambda$ two-$\rho$ terms is JIMWLK$\otimes$LO (\texttt{TwoRho\_from\_final.pdf}, \texttt{Evolution\_from\_final.pdf}). That sum is UV finite: in the Hamiltonian every kernel $K(u,v,t)$ multiplies $[U^\dagger(u)-U^\dagger(t)][U(v)-U(t)]$, which vanishes at $t\to u$ and $t\to v$. So the $1/(x-z)^2$ singularities of the single rows cancel in the sum. With all rows in dimensional regularization, their poles cancel in the same way.
\item The UV term plus the rest is your original expression. So the $\log\frac\vee\Lambda$ content of 2b used in those notes is unchanged.
\item The poles that multiply $\log\frac\vee{k^+}$, $\log(1-\frac\Lambda{k^+})$ and the powers of $\Lambda/k^+$ (and $\log(1-\frac{k^+}\vee)$ in G1-II) must cancel against the other rows in the same scheme. This has not been checked here.
\end{itemize}

%==========================================================================================
\section{The rest}\label{sec:rest}
Eq.~\eqref{eq:rest} is 2b minus the counterterm \eqref{eq:CT}, with the $p^+$ integral over $\Lambda<p^+<k^+-\Lambda$ as in your notes.
\begin{itemize}
\item \textbf{UV finite.} At $x\to z$ the $1/(x-z)^2$ terms cancel for every $p^+$: Eq.~\eqref{eq:limit} and property (i). What is left has zero angular average or is $O(1/|x-z|)$. So the $x$ integral converges at short distance, and the rest can be evaluated in two dimensions ($d_\perp=2$).

  Numerically, with a smooth test colour factor and the phase kept:
  \begin{itemize}
  \item $r^2\langle\text{rest}\rangle\to0$ at $p^+/k^+=0.02,\,0.37,\,0.95$;
  \item $\int_{r_0<|x-z|<0.3}d^2x$ of the rest converges as $r_0\to0$, while the same integral of 2b grows like $2\pi\log(1/r_0)$ (Table~\ref{tab:checks}).
  \end{itemize}
\item \textbf{Large $|x-z|$.} The rest is 2b there, since the counterterm falls like $|x-z|^{-3}$. As for 2b itself, the $x$ integral converges through the phase $e^{-ik\cdot(z-x)p^+/k^+}$.
\item \textbf{Rapidity logarithms.} The rest still has the $\frac{k^+}{p^+}$ and $\frac{k^+}{k^+-p^+}$ behaviour at the ends of the $p^+$ integral, from $B_4$ and $B_2$ and from the counterterm. So it still contains $\log\frac{k^+-\Lambda}{\Lambda}$ terms, now with UV-finite kernels. No plus prescription is applied; the $p^+$ integral is over $\Lambda<p^+<k^+-\Lambda$.
\item \textbf{The c.c.}\ is the complex conjugate of Eqs.~\eqref{eq:UV} and \eqref{eq:rest}; it reverses the order of the charges.
\end{itemize}

%==========================================================================================
\section{Checks}\label{sec:checks}
All numbers come from \texttt{part2b\_checks.py} (\texttt{part2b\_checks.json}). In the checks $z=0$, $y=(0.8,-0.5)$, $\mathbf k=(1.3,0.4)$. The colour factor $N_c[U(y)-U(x)]$ is replaced by a smooth test function $G(y)-G(x)$, and the counterterm gets $G(y)-G(z)$.
\begin{table}[H]
\centering\small
\begin{tabular}{@{}p{0.50\textwidth} p{0.45\textwidth}@{}}
check & result\\ \hline
short distance: $r^2\langle\text{integrand}\rangle\,\frac{(y-z)^2}{(y-z)^i}$ vs.\ $\xi(1-\xi)+\frac{\xi}{1-\xi}+\frac{1-\xi}{\xi}$; $d=2$ and $d=3$ ($d_\perp=d$); $\xi=0.1,0.37,0.8$; $r=10^{-4}$ & max relative difference $1.6\times10^{-8}$\\
Eq.~\eqref{eq:angle}: $r^2\langle K(z,z,x)-K(y,z,x)\rangle$ at $|x-z|/|y-z|=0.01, 0.32, 0.90, 1.10, 3.18, 31.80$ & 1.000000, 1.000000, 1.000000, 0.000000, 0.000000, 0.000000\\
$\xi=0.02$: $r^2|\langle{\rm 2b}\rangle|$ and $r^2|\langle\text{rest}\rangle|$ at $r=10^{-2},10^{-3},10^{-4}$ & 16.7232, 16.7240, 16.7240; $8.6\times10^{-4}$, $8.6\times10^{-6}$, $8.6\times10^{-8}$\\
$\xi=0.37$: $r^2|\langle{\rm 2b}\rangle|$ and $r^2|\langle\text{rest}\rangle|$ at $r=10^{-2},10^{-3},10^{-4}$ & 0.8604, 0.8604, 0.8604; $5.9\times10^{-5}$, $5.9\times10^{-7}$, $5.9\times10^{-9}$\\
$\xi=0.95$: $r^2|\langle{\rm 2b}\rangle|$ and $r^2|\langle\text{rest}\rangle|$ at $r=10^{-2},10^{-3},10^{-4}$ & 6.5127, 6.5137, 6.5137; $1.1\times10^{-3}$, $1.1\times10^{-5}$, $1.1\times10^{-7}$\\
$\xi=0.2$: ${\rm Re}\int_{r_0<|x-z|<0.3}d^2x$ (component $i=1$), $r_0=10^{-2},\dots,10^{-5}$: rest; 2b & 0.250416, 0.250689, 0.250692, 0.250692; $-31.89,\ -53.65,\ -75.41,\ -97.16$\\
$\xi=0.7$: ${\rm Re}\int_{r_0<|x-z|<0.3}d^2x$ (component $i=1$), $r_0=10^{-2},\dots,10^{-5}$: rest; 2b & 0.300763, 0.301098, 0.301102, 0.301102; $-21.36,\ -36.02,\ -50.68,\ -65.35$\\
same: 2b minus $\pi\log\frac{0.09}{r_0^2}\times$(coefficient of Eq.~\eqref{eq:limit}) minus rest & max $2.5\times10^{-14}$\\
$\xi=0.2$, $|x-z|=10,10^2,10^3$: $|x-z|^3\max|{\rm CT}|$; $|x-z|^2\max|{\rm 2b}|$; $|x-z|^2|\langle{\rm 2b}\rangle|$ without phase & 1.567, 1.432, 1.420; 5.186, 4.956, 4.920; 2.287, 2.303, 2.303\\
$\xi=0.7$, $|x-z|=10,10^2,10^3$: $|x-z|^3\max|{\rm CT}|$; $|x-z|^2\max|{\rm 2b}|$; $|x-z|^2|\langle{\rm 2b}\rangle|$ without phase & 1.056, 0.965, 0.957; 2.624, 3.736, 3.627; 1.165, 1.166, 1.166\\
$\xi=0.2$: ${\rm Re}\int_{1.5<|x-z|<R_{\max}}d^2x$ rest (phase kept), $R_{\max}=25,50,100,200$ & $6.785,\ 6.981,\ 7.001,\ 6.961$\\
$\xi=0.7$: ${\rm Re}\int_{1.5<|x-z|<R_{\max}}d^2x$ rest (phase kept), $R_{\max}=25,50,100,200$ & $-1.213,\ -1.171,\ -1.182,\ -1.181$\\
$d_\perp=2-\epsilon$ integrals minus $\pi[\frac2\epsilon+2\gamma_E+\log\frac{R^2\mu^2}{4}]$: $K$ form (Eq.~\eqref{eq:dimreg}), cut, smooth form (\texttt{sympy}) & 0; 0; 0\\
$K$-form formula at $d=3$ minus $\frac14\cdot\frac{4\pi}{R}$ (\texttt{sympy}) & 0\\
Eq.~\eqref{eq:xi}: exact integral minus the right-hand side at $\lambda=\frac13,\frac1{10},\frac1{1000}$ (\texttt{sympy}, 30 digits) & $5.4\times10^{-166}$\\
prefactor of the UV term$/\frac{1}{(2\pi)^3}\frac{g^4N_c}{16\pi^4}\frac{1}{k^+}$; $/\big(\frac{\alpha_sN_c}{2\pi}\times$LO prefactor$\big)$ & 1; 1\\
$\frac{\alpha_sN_c}{2\pi}(-\frac{11}{6})+\frac{\alpha_s\beta_0}{4\pi}$; $\beta_0$ coefficient G1-II(1) + 2b & 0; 0\\
\end{tabular}
\caption{Numerical and symbolic checks.}
\label{tab:checks}
\end{table}

\end{document}
